\section{Fixing the policy question before the search finishes}
\label{sim:prospective-partial-policy}

A manager who may pay an obligation or draw on a facility needs a rule for acting on the information available at each decision. The preceding search can establish some possible information histories while leaving others unexplored. To compare policies during that search, their meaning must remain fixed. If each newly discovered history changes what a policy is allowed to observe, an apparent reduction in the feasible set may instead reflect a change in the question.

The first requirement is therefore a declaration of possible messages before temporal search. In the finite example, start with the original registry of financial assertions and observation schemes. For each assertion-valued observation scheme, declare the first- and second-decision messages for every raw observation obtainable from the original financial assertions. Cross those financial and observation assertions conservatively. Some combinations may later prove temporally impossible; including them at this stage avoids selecting a favourable source history in advance.

The declaration binds the manager, decision dates, source records and meanings of pay and draw. A second-decision key also retains the manager's own earlier action, as in the preceding policy example. Internal source-head identifiers remain evidence of support and do not become messages to the manager. A retraction supplies no replacement assertion payload, but it does not justify removing every message associated with the original assertion before its timing has been examined. This is a declared finite vocabulary, not a finding that a real manager received every listed message.

\subsection{A fixed domain makes partial restrictions comparable}

Let $U$ denote the declared finite set of decision keys, including the decision stage and the required own-action memory. A total rule assigns either pay or draw to every key in $U$. Let $\Pi_U$ be the set of these rules. Assignments to keys that are never reached still belong to a total rule; they have no effect on that rule's realised path. Consequently, several total rules can represent the same policy on the reachable histories.

At search stage $t$, let $H_t$ contain only the selected financial and information histories supported by completed temporal branches. For a history $h$, write $V_h(\pi)$ when following rule $\pi$ selects no refused action in that history. This condition uses the complete finite action rows for that selected history, including first-decision refusal and second-decision refusal after the first action. The rules surviving the known restrictions are
\begin{equation}
 C_t=\{\pi\in\Pi_U:V_h(\pi)\text{ for every }h\in H_t\}.
 \label{eq:prospective-policy-restriction}
\end{equation}
These are sets of decision rules, not monetary balances or probability distributions. The relation is a local definition under a fixed message declaration and the supplied action semantics. Unresolved source outcomes are retained separately; they cannot supply a selected history or establish its financial constraints.

Resuming the same checked search retains the previously completed histories. Thus $H_t\subseteq H_{t+1}$. A rule satisfying every condition for the larger set also satisfies every condition for its subset, so
\begin{equation}
 C_{t+1}\subseteq C_t.
 \label{eq:prospective-policy-monotone}
\end{equation}
This inclusion explains the asymmetric conclusions. Once known constraints eliminate every rule, adding more histories cannot restore a rule under the same declaration. If some rules survive, a pending branch may still eliminate them. In the coarse/refined correction example, the known constraints already eliminate all rules after four of the 25 temporal nodes have been resolved. The baseline instead retains four total rules and establishes conditional feasibility only when the search completes with selected source outcomes throughout.

\subsection{Three reasons that a surviving rule may still be unresolved}

Unfinished search, unresolved source meaning and an inadequate domain allowance are different limitations. A two-node search cap retains 15 pending branches in the baseline. Those branches have not all been shown to contain possible histories, but they must remain represented until checked. The surviving rules therefore remain candidates.

A search can also finish while a required assertion remains unavailable, retracted or in conflict with another head. Unknown access leaves four rules surviving the known constraints, yet one source outcome remains unresolved. Retraction and conflicting-head cases likewise do not become feasible merely because all temporal branches have been processed. An inconsistent clock system is different again: it admits no temporal assignment. The 512 unconstrained rule assignments in that example do not establish a feasible policy for a possible company history.

Finally, the conservative domain may exceed the stated computational allowance. The finite profile permits at most 12 binary decision keys, hence at most 4,096 total rules. Renaming only the low message in one correction produces 12 keys and eight surviving total rules. Renaming all correction messages requires 18 keys, so the reference retains the original request and returns an unresolved domain budget without truncating the keys or enumerating a preferred subset. The extra rule count in the smaller example is not evidence of better information or better economic performance.

The checked partial calculation binds each result to the original request, declaration, domain, methods and temporal frontier. It retains the supporting branches, unresolved source outcomes and pending temporal domains. Changing a message, source, actor or domain cannot reuse the old result as the answer to the new question. A future extension rule may justify some reuse, but that relation must itself be stated and checked.\footnote{The original-registry domain comparison is retained in \path{data/accounting-master/prospective-policy-domain-20261027/}; its 10 cases include explicit overflow and conflicting-head cases. The partial comparison in \path{data/accounting-master/prospective-partial-policy-20261027/} checks 12 cases at 110 frontiers, with independent full-rule enumeration and 31 refusal controls. The candidate takes no complete-family comparison as an input. Exact sets and endpoint projections, rather than the counts alone, support these finite conclusions.}

This construction joins a fixed decision problem to an unfinished search without treating uncertainty as missing arithmetic. Longer horizons, additional actions, changing declarations and broader observation languages still require their own comparisons.

